% Mechanistic Validity — NEMI 2026 poster, v3
% Built from the actual paper. All tier names, criteria, and numbers verified.
%
%   lualatex mechval_poster_grok_v3.tex

\documentclass[25pt, a0paper, portrait]{tikzposter}

\usepackage{amsmath,amssymb}
\usepackage{array}
\usepackage{booktabs}
\usepackage{graphicx}
\usepackage{enumitem}
\usepackage{microtype}

% ── Color palette ───────────────────────────────────────
\definecolor{confirmed}{HTML}{27AE60}
\definecolor{partial}{HTML}{2980B9}
\definecolor{inconclusive}{HTML}{F39C12}
\definecolor{untested}{HTML}{95A5A6}
\definecolor{disconfirmed}{HTML}{E74C3C}
\definecolor{darktext}{HTML}{2C3E50}
\definecolor{blockbg}{HTML}{FFFFFF}
\definecolor{tierProposed}{HTML}{95A5A6}
\definecolor{tierCausal}{HTML}{F39C12}
\definecolor{tierSupported}{HTML}{2980B9}
\definecolor{tierTriangulated}{HTML}{27AE60}
\definecolor{tierValidated}{HTML}{8E44AD}
\definecolor{tierUndetermined}{HTML}{E67E22}
\definecolor{tierDisconfirmed}{HTML}{E74C3C}

% ── Theme ───────────────────────────────────────────────
\usetheme{Default}

\definecolorstyle{MechValColors}{
  \definecolor{colorOne}{HTML}{2C3E50}
  \definecolor{colorTwo}{HTML}{3498DB}
  \definecolor{colorThree}{HTML}{ECF0F1}
}{
  \colorlet{backgroundcolor}{colorThree}
  \colorlet{framecolor}{colorOne}
  \colorlet{titlefgcolor}{white}
  \colorlet{titlebgcolor}{colorOne}
  \colorlet{blocktitlebgcolor}{colorTwo}
  \colorlet{blocktitlefgcolor}{white}
  \colorlet{blockbodybgcolor}{blockbg}
  \colorlet{blockbodyfgcolor}{darktext}
  \colorlet{notefgcolor}{darktext}
  \colorlet{notebgcolor}{inconclusive!15}
  \colorlet{notefrcolor}{inconclusive}
}
\usecolorstyle{MechValColors}

\defineblockstyle{MechValBlock}{
  titlewidthscale=1.0, bodywidthscale=1.0, titleleft,
  titleoffsetx=0pt, titleoffsety=0pt,
  bodyoffsetx=0pt, bodyoffsety=0pt, bodyverticalshift=0pt,
  roundedcorners=12, linewidth=2pt, titleinnersep=10pt, bodyinnersep=15pt
}{
  \draw[rounded corners=\blockroundedcorners, color=blocktitlebgcolor,
        fill=blockbodybgcolor, line width=\blocklinewidth]
    (blockbody.south west) rectangle (blocktitle.north east);
  \ifBlockHasTitle
    \draw[rounded corners=\blockroundedcorners,
          fill=blocktitlebgcolor, color=blocktitlebgcolor,
          line width=\blocklinewidth]
      (blocktitle.south west) rectangle (blocktitle.north east);
  \fi
}
\useblockstyle{MechValBlock}

\definetitlestyle{MechValTitle}{
  width=\paperwidth, roundedcorners=0, linewidth=0pt, innersep=20pt,
  titletotopverticalspace=15mm, titletoblockverticalspace=25mm
}{
  \draw[fill=titlebgcolor, color=titlebgcolor]
    (\titleposleft, \titleposbottom)
    rectangle (\titleposright, \titlepostop);
}
\usetitlestyle{MechValTitle}

% verdict bar helper: \vbar{color}{width mm}
\newcommand{\vbar}[2]{\textcolor{#1}{\rule[0.6mm]{#2mm}{5mm}}}

\title{%
  \parbox{0.92\linewidth}{\centering\bfseries
    Mechanistic Validity\\[6pt]
    A Validity Theory, Research Methodology, and Evidence Standard\\[3pt]
    for Mechanistic Interpretability}}
\author{Elliot Tower\quad\texttt{elliot@elliottower.ai}}
\institute{}

\begin{document}

\maketitle[width=\paperwidth]

% ════════════════════════════════════════════════════════
% ROW 1
% ════════════════════════════════════════════════════════
\begin{columns}
\column{0.33}

\block{The Problem}{
  A mechanistic claim is only as strong as its weakest validity
  dimension. Five validity types---construct, measurement, internal,
  external, and interpretive---form a dependency chain: construct
  validity is a ceiling on all others.

  \vspace{6mm}
  No amount of causal evidence rescues a claim whose target concept is
  incoherent. No measurement reliability compensates for measuring
  the wrong thing. A causal result established in one setting cannot
  automatically support broader generalization.

  \vspace{6mm}
  \textbf{The field's dominant failure mode:} establishing \emph{that}
  a component matters while leaving open \emph{what} it is.
}

\column{0.34}

\block{The Framework}{
  \textbf{36 criteria across 5 validity types.} Each claim receives a
  structured validity profile. The tier is bounded by the weakest
  dimension, not the strongest.

  \vspace{5mm}
  \begin{center}
  \begin{tabular}{lcp{0.52\linewidth}}
    \toprule
    \textbf{Type} & \textbf{Criteria} & \textbf{Core question} \\
    \midrule
    Construct    & C1--C6  & Is the circuit a natural kind? \\
    Measurement  & M1--M7  & Does the metric measure what it claims? \\
    Internal     & I1--I12 & Is the causal claim warranted? \\
    External     & E1--E6  & Does the result generalize? \\
    Interpretive & V1--V5  & Is the interpretation justified? \\
    \bottomrule
  \end{tabular}
  \end{center}

  \vspace{4mm}
  The criteria form a dependency chain: the target mechanism must be
  coherent before it can be measured; the measurement must be
  trustworthy before it can support causal inference.
}

\column{0.33}

\block{The Tier System}{
  Five verdict tiers, each subsuming all requirements below it.

  \vspace{5mm}
  \begin{enumerate}[leftmargin=*, itemsep=3mm]
    \item[\textcolor{tierProposed}{\textbullet}] \textbf{Proposed:}
      structural or representational evidence only
    \item[\textcolor{tierCausal}{\textbullet}] \textbf{Causally Suggestive:}
      necessity shown, sufficiency not established
    \item[\textcolor{tierSupported}{\textbullet}] \textbf{Mechanistically Supported:}
      necessity + sufficiency with consistent methods
    \item[\textcolor{tierTriangulated}{\textbullet}] \textbf{Triangulated:}
      converging evidence from independent methods
    \item[\textcolor{tierValidated}{\textbullet}] \textbf{Validated:}
      all five validity types addressed
  \end{enumerate}

  \vspace{4mm}
  Three diagnostic labels replace the tier when a claim falls outside
  the hierarchy: \textbf{Underdetermined} (evidence fits multiple
  mechanisms), \textbf{Insufficient} (cannot assess), and
  \textbf{Disconfirmed} (fails decisively on a required criterion).
}

\end{columns}

% ════════════════════════════════════════════════════════
% ROW 2 — CASE STUDIES + IOI AUDIT
% ════════════════════════════════════════════════════════
\begin{columns}
\column{0.50}

\block{Sixteen Claims from Fifteen Papers}{
  {\small
  \begin{tabular}{@{}>{\raggedright\arraybackslash}p{3.0cm}>{\raggedright\arraybackslash}p{6.5cm}l@{}}
    \toprule
    \textbf{Verdict} & \textbf{Claim} & \textbf{Scope} \\
    \midrule
    \textcolor{tierTriangulated}{\textbf{Triangulated}}
      & Induction heads: token copying & 1L--70B \\
    \midrule
    \textcolor{tierSupported}{\textbf{Mech.\ Supported}}
      & Fourier mechanism (mod.\ addition) & 1L ReLU \\
      & Greater-Than circuit & GPT-2 sm \\
      & Copy Suppression (L10H7) & GPT-2 sm \\
      & Superposition (toy networks) & 4 ReLU nets \\
      & Steering vector (refusal) & 13 models \\
      & Global Workspace / J-space & 4 Claude \\
      & Successor Heads & Pythia-1.4B \\
    \midrule
    \textcolor{tierCausal}{\textbf{Causally Suggestive}}
      & Othello World Model & Othello-GPT \\
      & Docstring Circuit & 4L attn-only \\
      & Gender Bias Circuits & GPT-2 + 5 \\
    \midrule
    \textcolor{tierProposed}{\textbf{Proposed}}
      & Probing Classifiers & review \\
    \midrule
    \textcolor{tierUndetermined}{\textbf{Undetermined}}
      & IOI Circuit & GPT-2 sm \\
    \midrule
    \textcolor{tierDisconfirmed}{\textbf{Disconfirmed}}
      & Knowledge Neurons & BERT-base \\
      & SAE features as canonical units & Pythia \\
      & Induction heads: general ICL & $>$1B \\
    \bottomrule
  \end{tabular}
  }

  \vspace{4mm}
  No claim reaches \textbf{Validated}. The highest tier achieved is
  Triangulated, by induction heads for token copying---and only the
  narrow claim, not the general one. Same paper, two claims, two
  tiers.
}

\column{0.50}

\block{IOI Circuit: Full 36-Criterion Audit}{
  {\small
  \begin{tabular}{@{}clcl@{}}
    \toprule
    & \textbf{Construct} & & \textbf{Internal} \\
    \midrule
    \textcolor{confirmed}{\textbf{C}} & C1 Falsifiability
      & \textcolor{partial}{\textbf{PC}} & I1 Necessity \\
    \textcolor{partial}{\textbf{PC}} & C2 Structural plausibility
      & \textcolor{partial}{\textbf{PC}} & I2 Sufficiency \\
    \textcolor{partial}{\textbf{PC}} & C3 Convergent validity
      & \textcolor{inconclusive}{\textbf{I}} & I3 Minimality \\
    \textcolor{inconclusive}{\textbf{I}} & C4 Discriminant validity
      & \textcolor{inconclusive}{\textbf{I}} & I4 Specificity \\
    \textcolor{partial}{\textbf{PC}} & C5 Nomological validity
      & \textcolor{inconclusive}{\textbf{I}} & I5 Rival exclusion \\
    \textcolor{untested}{\textbf{U}} & C6 Complementation
      & \textcolor{untested}{\textbf{U}} & I6 Double dissociation \\
    \midrule
    & \textbf{Measurement} & \textcolor{partial}{\textbf{PC}} & I7 Confound control \\
    \cmidrule{1-2}
    \textcolor{inconclusive}{\textbf{I}} & M1 Reliability
      & \textcolor{untested}{\textbf{U}} & I8 Confound sensitivity \\
    \textcolor{partial}{\textbf{PC}} & M2 Baseline separation
      & \textcolor{partial}{\textbf{PC}} & I9 Epistatic interaction \\
    \textcolor{disconfirmed}{\textbf{D}} & M3 Stability
      & \textcolor{partial}{\textbf{PC}} & I10 Rescue reversibility \\
    \textcolor{partial}{\textbf{PC}} & M4 Calibration
      & \textcolor{partial}{\textbf{PC}} & I11 Onset coupling \\
    \textcolor{untested}{\textbf{U}} & M5 Sensitivity
      & \textcolor{partial}{\textbf{PC}} & I12 Offset coupling \\
    \textcolor{disconfirmed}{\textbf{D}} & M6 Invariance & & \\
    \textcolor{untested}{\textbf{U}} & M7 Selection correction & & \\
    \midrule
    & \textbf{External} & & \textbf{Interpretive} \\
    \midrule
    \textcolor{inconclusive}{\textbf{I}} & E1 Intervention reach
      & \textcolor{partial}{\textbf{PC}} & V1 Level declaration \\
    \textcolor{disconfirmed}{\textbf{D}} & E2 Prompt generalization
      & \textcolor{partial}{\textbf{PC}} & V2 Level-evidence match \\
    \textcolor{partial}{\textbf{PC}} & E3 Cross-task
      & \textcolor{inconclusive}{\textbf{I}} & V3 Alternative level \\
    \textcolor{partial}{\textbf{PC}} & E4 Cross-model
      & \textcolor{partial}{\textbf{PC}} & V4 Unlicensed labeling \\
    \textcolor{partial}{\textbf{PC}} & E5 Graded response
      & \textcolor{confirmed}{\textbf{C}} & V5 Scope declaration \\
    \textcolor{confirmed}{\textbf{C}} & E6 Novel prediction & & \\
    \bottomrule
  \end{tabular}
  }

  \vspace{4mm}
  \begin{center}
  \begin{tabular}{lcc}
    \toprule
    \textbf{Verdict} & \textbf{Count} & \\
    \midrule
    \textcolor{confirmed}{\textbf{Confirmed}} & 3 &
      \vbar{confirmed}{18} \\
    \textcolor{partial}{\textbf{Partially confirmed}} & 18 &
      \vbar{partial}{108} \\
    \textcolor{inconclusive}{\textbf{Inconclusive}} & 7 &
      \vbar{inconclusive}{42} \\
    \textcolor{untested}{\textbf{Untested}} & 5 &
      \vbar{untested}{30} \\
    \textcolor{disconfirmed}{\textbf{Disconfirmed}} & 3 &
      \vbar{disconfirmed}{18} \\
    \bottomrule
  \end{tabular}
  \end{center}

  \vspace{3mm}
  \textbf{Verdict: Causally Suggestive.} Scored at the claim's most
  lenient reading, where the rival head sets that reproduce the behavior
  leave it standing. Method disagreement caps it: the headline 87\%
  faithfulness spans below 0\% to over 100\% across six methodological
  choices. The field's most studied mechanism, read as generously as the
  evidence allows, still stops two tiers short.
}

\end{columns}

% ════════════════════════════════════════════════════════
% ROW 3 — SOURCE DISCIPLINES
% ════════════════════════════════════════════════════════
\begin{columns}
\column{0.33}

\block{From Philosophy of Science}{
  Popper's falsifiability, Mayo's severe testing, Glennan's new
  mechanical philosophy. A circuit discovered by activation patching
  and evaluated by activation patching has only been shown to pass the
  test it was built to pass.

  \vspace{4mm}
  \textbf{Corroboration requires a genuinely risky test the circuit
  was not optimized for.}
}

\column{0.34}

\block{From Measurement Theory}{
  Cronbach and Meehl's construct validity (1955): a measurement must
  converge with measures of the same construct and diverge from
  measures of different constructs.

  \vspace{4mm}
  \textbf{Convergent without discriminant is not enough}---the same
  heads appear in most circuits, and no published claim tests what its
  components do \emph{not} do.
}

\column{0.33}

\block{From Causal Inference}{
  Rubin's potential outcomes, Rosenbaum's sensitivity analysis,
  E-values for unmeasured confounders.

  \vspace{4mm}
  \textbf{Ablation is an intervention, and every intervention needs a
  sensitivity analysis.} The field borrows the interventional
  vocabulary; it should borrow the statistical discipline too.
}

\end{columns}

% ════════════════════════════════════════════════════════
% ROW 4 — SOURCE DISCIPLINES (continued) + TAKEAWAY
% ════════════════════════════════════════════════════════
\begin{columns}
\column{0.33}

\block{From Neuroscience}{
  Craver's mutual manipulability, constitutive relevance criteria.
  A mechanism must be located at a level and testable by intervention
  at that level.

  \vspace{4mm}
  \textbf{The level at which a claim is made determines what evidence
  it needs}---weight-level claims need weight-level tests.
}

\column{0.34}

\block{From Pharmacology}{
  Dose-response curves, receptor theory, the ORBITA trial.

  \vspace{4mm}
  \textbf{Graded interventions distinguish specific from nonspecific
  effects.} A circuit ablation is a maximal dose; the dose-response
  curve between zero and complete ablation carries more information
  than the endpoint alone.
}

\column{0.33}

\block{From Genetics}{
  Epistasis mapping, rescue experiments, sensitivity analysis,
  complementation tests (Benzer's cis-trans procedure).

  \vspace{4mm}
  \textbf{Interaction structure, functional rescue, construct
  boundaries.} These fill gaps no existing MI technique covers.

  \vspace{8mm}
  \begin{center}
  {\large\textbf{Takeaway}}

  \vspace{4mm}
  A mechanistic claim is warranted only to the\\
  extent that every required link between concept,\\
  measurement, causal evidence, scope, and\\
  interpretation is supported.

  \vspace{6mm}
  \texttt{elliot@elliottower.ai}

  \vspace{2mm}
  {\small Paper:}
  \texttt{doi.org/10.5281/zenodo.20478480}
  \end{center}
}

\end{columns}

\end{document}
